% Options for packages loaded elsewhere
\PassOptionsToPackage{unicode}{hyperref}
\PassOptionsToPackage{hyphens}{url}
%
\documentclass[
  UTF8]{ctexart}
\usepackage{amsmath,amssymb}
\usepackage{iftex}
\ifPDFTeX
  \usepackage[T1]{fontenc}
  \usepackage[utf8]{inputenc}
  \usepackage{textcomp} % provide euro and other symbols
\else % if luatex or xetex
  \usepackage{unicode-math} % this also loads fontspec
  \defaultfontfeatures{Scale=MatchLowercase}
  \defaultfontfeatures[\rmfamily]{Ligatures=TeX,Scale=1}
\fi
\usepackage{lmodern}
\ifPDFTeX\else
  % xetex/luatex font selection
\fi
% Use upquote if available, for straight quotes in verbatim environments
\IfFileExists{upquote.sty}{\usepackage{upquote}}{}
\IfFileExists{microtype.sty}{% use microtype if available
  \usepackage[]{microtype}
  \UseMicrotypeSet[protrusion]{basicmath} % disable protrusion for tt fonts
}{}
\makeatletter
\@ifundefined{KOMAClassName}{% if non-KOMA class
  \IfFileExists{parskip.sty}{%
    \usepackage{parskip}
  }{% else
    \setlength{\parindent}{0pt}
    \setlength{\parskip}{6pt plus 2pt minus 1pt}}
}{% if KOMA class
  \KOMAoptions{parskip=half}}
\makeatother
\usepackage{xcolor}
\usepackage[margin=2cm]{geometry}
\usepackage{longtable,booktabs,array}
\usepackage{calc} % for calculating minipage widths
% Correct order of tables after \paragraph or \subparagraph
\usepackage{etoolbox}
\makeatletter
\patchcmd\longtable{\par}{\if@noskipsec\mbox{}\fi\par}{}{}
\makeatother
% Allow footnotes in longtable head/foot
\IfFileExists{footnotehyper.sty}{\usepackage{footnotehyper}}{\usepackage{footnote}}
\makesavenoteenv{longtable}
\setlength{\emergencystretch}{3em} % prevent overfull lines
\providecommand{\tightlist}{%
  \setlength{\itemsep}{0pt}\setlength{\parskip}{0pt}}
\setcounter{secnumdepth}{-\maxdimen} % remove section numbering
\ifLuaTeX
  \usepackage{selnolig}  % disable illegal ligatures
\fi
\usepackage{bookmark}
\IfFileExists{xurl.sty}{\usepackage{xurl}}{} % add URL line breaks if available
\urlstyle{same}
\hypersetup{
  hidelinks,
  pdfcreator={LaTeX via pandoc}}

\author{}
\date{}

\begin{document}

\textbf{1．分式什么时候等于零}

\[
\boxed{
\frac{N(x)}{D(x)}=0
\iff
N(x)=0\quad\text{且分母有定义、非零}.
}
\]

做题步骤：\textbf{令分子为零 → 检查候选解是否在原式定义域内。}

\begin{itemize}
\tightlist
\item
  分母再大，只要有限，就不能把非零分子除成恰好 \(0\)。
\item
  \textbf{趋近于 \(0\)} 与 \textbf{等于 \(0\)} 要区分。
\end{itemize}

今天的例子： \[
f'(x)=\frac{a-x}{(x^2-2x+a)^{3/2}}.
\] 候选零点是 \(x=a\)，但需要满足 \[
a^2-a>0
\iff
\boxed{a<0\text{ 或 }a>1}.
\]

\begin{center}\rule{0.5\linewidth}{0.5pt}\end{center}

\textbf{2．高次韦达定理}

设 \[
a_nx^n+a_{n-1}x^{n-1}+\cdots+a_0=0,\qquad a_n\ne0,
\] 根为 \(r_1,\ldots,r_n\)。定义 \(e_k\) 为：\textbf{选出 \(k\)
个不同编号的根，相乘后对所有选法求和}，共 \(\binom nk\) 项。

\[
\boxed{e_k=(-1)^k\frac{a_{n-k}}{a_n}}.
\]

具体为： \[
\begin{aligned}
e_1&=\sum_i r_i=-\frac{a_{n-1}}{a_n},\\
e_2&=\sum_{i<j}r_ir_j=\frac{a_{n-2}}{a_n},\\
e_3&=\sum_{i<j<k}r_ir_jr_k=-\frac{a_{n-3}}{a_n},\\
e_n&=\prod_i r_i=(-1)^n\frac{a_0}{a_n}.
\end{aligned}
\]

\textbf{重点：符号交替；都除以最高次项系数；每组选法只算一次。}
允许复根，重根按重数计入。

常用变形： \[
\boxed{\sum_i r_i^2=e_1^2-2e_2},
\] \[
\boxed{\sum_i r_i^3=e_1^3-3e_1e_2+3e_3}.
\] 所有根非零时： \[
\boxed{\sum_i\frac1{r_i}=\frac{e_{n-1}}{e_n}=-\frac{a_1}{a_0}}.
\]

原理：展开 \(a_n\prod_i(x-r_i)\)，比较系数。

\begin{center}\rule{0.5\linewidth}{0.5pt}\end{center}

\textbf{3．三次方程：等差根、等比根与代入技巧}

对 \[
f(x)=x^3-px^2+qx-r,\qquad p\ne0,\ r\ne0,
\] 韦达关系为： \[
\boxed{\alpha+\beta+\gamma=p,\quad
\alpha\beta+\beta\gamma+\gamma\alpha=q,\quad
\alpha\beta\gamma=r}.
\]

\begin{longtable}[]{@{}
  >{\raggedright\arraybackslash}p{(\columnwidth - 6\tabcolsep) * \real{0.2500}}
  >{\raggedright\arraybackslash}p{(\columnwidth - 6\tabcolsep) * \real{0.2500}}
  >{\raggedright\arraybackslash}p{(\columnwidth - 6\tabcolsep) * \real{0.2500}}
  >{\raggedright\arraybackslash}p{(\columnwidth - 6\tabcolsep) * \real{0.2500}}@{}}
\toprule\noalign{}
\begin{minipage}[b]{\linewidth}\raggedright
根的结构
\end{minipage} & \begin{minipage}[b]{\linewidth}\raggedright
推荐设根
\end{minipage} & \begin{minipage}[b]{\linewidth}\raggedright
中间根
\end{minipage} & \begin{minipage}[b]{\linewidth}\raggedright
系数充要条件
\end{minipage} \\
\midrule\noalign{}
\endhead
\bottomrule\noalign{}
\endlastfoot
等比 & \(a/k,\ a,\ ak\) & \(\displaystyle a=q/p\) &
\(\boxed{q^3=p^3r}\) \\
等差 & \(a-d,\ a,\ a+d\) & \(\displaystyle a=p/3\) &
\(\boxed{2p^3-9pq+27r=0}\) \\
\end{longtable}

\textbf{正向：把假设的根代入韦达，消掉辅助参数。}

等比： \[
p=a(k^{-1}+1+k),\quad
q=a^2(k^{-1}+1+k)=ap,\quad
r=a^3.
\]

等差： \[
p=3a,\qquad q=3a^2-d^2,\qquad r=a(a^2-d^2).
\]

\textbf{反向：先验证中间根，再利用另外两根的和或积。}

可直接使用： \[
\boxed{f(q/p)=\frac{q^3}{p^3}-r},
\qquad
\boxed{f(p/3)=-\frac{2p^3}{27}+\frac{pq}{3}-r}.
\]

找到中间根 \(a\) 后，设另外两根为 \(\beta,\gamma\)：

\begin{itemize}
\tightlist
\item
  \textbf{等比：}证明 \(\beta\gamma=a^2\)，则 \[
  \frac a\beta=\frac\gamma a.
  \]
\item
  \textbf{等差：}证明 \(\beta+\gamma=2a\)，则另外两根关于 \(a\)
  对称，可写为 \(a-d,a+d\)。
\end{itemize}

本题允许复根；若要求三个实根成等差，还需 \(p^2\ge3q\)。

\textbf{本题要补的解题习惯：}

\begin{longtable}[]{@{}ll@{}}
\toprule\noalign{}
题目线索 & 第一反应 \\
\midrule\noalign{}
\endhead
\bottomrule\noalign{}
\endlastfoot
给出根的形式 & \textbf{代入韦达}，或写出因式分解后比较系数 \\
要证明指定的 \(t\) 是根 & \textbf{直接计算 \(f(t)\)} \\
三根成等差／等比 & \textbf{围绕中间根对称设根} \\
前一问推出某个中间根 & 后一问反向证明时，\textbf{优先代入这个候选根} \\
\end{longtable}

\begin{center}\rule{0.5\linewidth}{0.5pt}\end{center}

\textbf{4．充分、必要与充要条件}

固定：\(C\) 是``条件''，\(R\) 是``结果''。

\begin{longtable}[]{@{}lll@{}}
\toprule\noalign{}
类型 & 含义 & 要证明的方向 \\
\midrule\noalign{}
\endhead
\bottomrule\noalign{}
\endlastfoot
充分条件 & \textbf{有它就够了}，能保证结果 & \(C\Rightarrow R\) \\
必要条件 & \textbf{没有它不行}，结果成立必须满足它 &
\(R\Rightarrow C\) \\
充要条件 & 两者互相推出 & \(C\Longleftrightarrow R\) \\
\end{longtable}

例如，针对``能被 \(6\) 整除''：

\begin{itemize}
\tightlist
\item
  是偶数：\textbf{必要但不充分}，反例为 \(4\)。
\item
  能被 \(12\) 整除：\textbf{充分但不必要}，反例为 \(6\)。
\item
  同时能被 \(2\) 和 \(3\) 整除：\textbf{充要}。
\end{itemize}

证明充要条件时：

\begin{enumerate}
\def\labelenumi{\arabic{enumi}.}
\tightlist
\item
  \textbf{必要性：假设结果成立，推出条件。}
\item
  \textbf{充分性：假设条件成立，推出结果。}
\end{enumerate}

\textbf{只从特殊根的形式推出系数关系，只完成了必要性；还必须反向证明。}

英文方向： \[
\boxed{P\text{ only if }Q\quad\Longleftrightarrow\quad P\Rightarrow Q}.
\] 因此 ``The derivative is zero only when \(x=a\)''
只说明：\textbf{导数为零必须有
\(x=a\)}；是否真的在这里为零，还要检查定义域。

\end{document}
